% hangulfont.tex -- the Hangul typeface choices for GWEB's Korean back end.
%
% kotexgweb.tex reads this file (\input hangulfont). Everything that names a
% Korean typeface lives here, so switching to a different Korean font -- say from
% Noto to another CJK family -- is a change to this one file; kotexgweb.tex need
% not be touched. It defines the body fonts \tenmj ... \twelvemj, the math-hangul
% fonts, the title companions \hangultitlefont/\hantitlefont, and \sectionfont.
%
% luatexko provides Hangul shaping (and automatic fallback, so Korean appears
% even through gwebmac's CM fonts). The plain-TeX hangul setup is borrowed from
% luatexko-doc.tex: a serif body (Noto Serif), a sans group/title (Noto Sans),
% and a monospace face (Noto Sans Mono) paired with the matching CM fonts.
%
% Sizes.  Hangul fills nearly its whole em box while CM's lowercase reaches only
% an unusually small x-height (0.431em, against 0.511em for the Latin that Noto
% itself ships), so Hangul set at the nominal size looks too big next to CM.
% Matching the ink each script actually spans -- CM's ascender-to-descender is
% 0.889em, an average Hangul syllable 0.924em -- gives 0.962 of the Latin size;
% against cmtt, whose span is only 0.833em, the monospace face wants 0.93.  So
% each Hangul font below is (its CM partner's design size) times that ratio.  CM
% keeps the same proportions at every design size (x-height is 0.431em in cmr10,
% cmr9 and cmr8 alike), which is why a constant ratio is right and the fixed
% -0.2pt offset used before was not: it drifted from .980 at ten point to .971
% at seven.  The interhangul/interlatincjk/charraise settings are all in em and
% therefore follow the size on their own; they are not restated per size.
\input luatexko.sty
% Deferred loading.  \font reads the font the moment it runs, so without this
% every face named below is read at startup whether or not the document ever
% selects it -- some two seconds of a run, most of it for faces nobody asked
% for.  So we hold the specification as text and run the real \font on the
% switch's first use; a document pays only for what it sets.  \globaldefs makes
% that first load outlive whatever group we happen to be in (the first use is
% usually inside one, as in {\tenbd ...}), so each face is still read once per
% run.  The six math fonts below stay eager: \setmathhangulfonts wants real
% font identifiers to hand to \textfont, and cannot take a macro.
\def\lazyhangulfont#1#2{\protected\gdef#1{\lazyhfload#1{#2}}}
\def\lazyhfload#1#2{\begingroup\globaldefs=1 \sethangulfont#1#2\endgroup #1}
\def\lazyplainfont#1#2{\protected\gdef#1{\begingroup\globaldefs=1 \font#1#2\endgroup #1}}

\lazyhangulfont\tenmj{="Noto Serif CJK KR Light:%
  language=KOR;%
  script=hang;%
  interhangul=-0.04em;%
  interlatincjk=.125em;%
  charraise=-0.06em;%
  +compresspunctuations;%
  expansion=default;" at 9.62pt}
\lazyhangulfont\tentz{="Noto Sans Mono CJK KR:%
  language=KOR;%
  script=hang;%
  interhangul=-0.04em;%
  interlatincjk=.125em;%
  charraise=-0.03em;%
  +compresspunctuations;%
  expansion=default;" at 9.3pt}
\lazyhangulfont\tenbd{="Noto Serif CJK KR Bold:%
  language=KOR;%
  script=hang;%
  interhangul=-0.04em;%
  interlatincjk=.125em;%
  charraise=-0.06em;%
  +compresspunctuations;%
  expansion=default;" at 9.62pt}
\lazyhangulfont\tengt{="Noto Sans CJK KR Light:%
  language=KOR;%
  script=hang;%
  interhangul=-0.04em;%
  interlatincjk=.125em;%
  charraise=-0.06em;%
  +compresspunctuations;%
  expansion=default" at 9.62pt}
\lazyhangulfont\tengtd{="Noto Sans CJK KR DemiLight:%
  language=KOR;%
  script=hang;%
  interhangul=-0.04em;%
  interlatincjk=.125em;%
  charraise=-0.06em;%
  +compresspunctuations;%
  expansion=default" at 9.62pt}

\lazyhangulfont\ninemj{="Noto Serif CJK KR Light:%
  language=KOR;%
  script=hang;%
  interhangul=-0.04em;%
  interlatincjk=.125em;%
  charraise=-0.06em;%
  +compresspunctuations;%
  expansion=default" at 8.66pt}
\lazyhangulfont\ninetz{="Noto Sans Mono CJK KR:%
  language=KOR;%
  script=hang;%
  interhangul=-0.04em;%
  interlatincjk=.125em;%
  charraise=-0.03em;%
  +compresspunctuations;%
  expansion=default;" at 8.37pt}
\lazyhangulfont\ninebd{="Noto Serif CJK KR Bold:%
  language=KOR;%
  script=hang;%
  interhangul=-0.04em;%
  interlatincjk=.125em;%
  charraise=-0.06em;%
  +compresspunctuations;%
  expansion=default" at 8.66pt}
\lazyhangulfont\ninegt{="Noto Sans CJK KR Light:%
  language=KOR;%
  script=hang;%
  interhangul=-0.04em;%
  interlatincjk=.125em;%
  charraise=-0.06em;%
  +compresspunctuations;%
  expansion=default" at 8.66pt}
\lazyhangulfont\ninegtd{="Noto Sans CJK KR DemiLight:%
  language=KOR;%
  script=hang;%
  interhangul=-0.04em;%
  interlatincjk=.125em;%
  charraise=-0.06em;%
  +compresspunctuations;%
  expansion=default" at 8.66pt}

\lazyhangulfont\eightmj{="Noto Serif CJK KR Light:%
  language=KOR;%
  script=hang;%
  interhangul=-0.04em;%
  interlatincjk=.125em;%
  charraise=-0.06em;%
  +compresspunctuations;%
  expansion=default" at 7.7pt}
\lazyhangulfont\eighttz{="Noto Sans Mono CJK KR:%
  language=KOR;%
  script=hang;%
  interhangul=-0.04em;%
  interlatincjk=.125em;%
  charraise=-0.03em;%
  +compresspunctuations;%
  expansion=default" at 7.44pt}
\lazyhangulfont\eightbd{="Noto Serif CJK KR Bold:%
  language=KOR;%
  script=hang;%
  interhangul=-0.04em;%
  interlatincjk=.125em;%
  charraise=-0.06em;%
  +compresspunctuations;%
  expansion=default" at 7.7pt}
\lazyhangulfont\eightgt{="Noto Sans CJK KR Light:%
  language=KOR;%
  script=hang;%
  interhangul=-0.04em;%
  interlatincjk=.125em;%
  charraise=-0.06em;%
  +compresspunctuations;%
  expansion=default" at 7.7pt}
\lazyhangulfont\eightgtd{="Noto Sans CJK KR DemiLight:%
  language=KOR;%
  script=hang;%
  interhangul=-0.04em;%
  interlatincjk=.125em;%
  charraise=-0.06em;%
  +compresspunctuations;%
  expansion=default" at 7.7pt}

\lazyhangulfont\sevenmj{="Noto Serif CJK KR Light:%
  language=KOR;%
  script=hang;%
  interhangul=-0.04em;%
  interlatincjk=.125em;%
  charraise=-0.06em;%
  +compresspunctuations;%
  expansion=default" at 6.73pt}

\lazyhangulfont\twelvemj{="Noto Serif CJK KR:%
  language=KOR;%
  script=hang;%
  interhangul=-0.04em;%
  interlatincjk=.125em;%
  charraise=-0.06em;%
  +compresspunctuations;%
  expansion=default" at 11.54pt}

% hangul in math mode
\font\texthangul="Noto Serif CJK KR Regular" at 9.62pt
\font\texthangulnine="Noto Serif CJK KR Regular" at 8.66pt
\font\texthanguleight="Noto Serif CJK KR Regular" at 7.7pt
\font\scripthangul="Noto Serif CJK KR Regular" at 6.73pt
\font\scripthangulsix="Noto Serif CJK KR Regular" at 5.77pt
\font\scriptscripthangul="Noto Serif CJK KR Medium" at 4.81pt

% Title companions for the contents-page title (\con) and the typewriter limbo
% title: kotexgweb.tex pairs these Hangul faces with gwebmac's Latin \titlefont
% and \ttitlefont. The ratios follow the body fonts', 0.962 of the serif partner
% and 0.93 of the typewriter one.
\lazyhangulfont\hangultitlefont{="Noto Serif CJK KR Medium:%
  language=KOR;script=hang;interhangul=-0.04em;interlatincjk=.125em;%
  charraise=-0.06em;+compresspunctuations" at 13.96pt}

\lazyhangulfont\hantitlefont{="Noto Sans Mono CJK KR:%
  language=KOR;script=hang;interhangul=-0.04em;interlatincjk=.125em;%
  charraise=-0.06em;+compresspunctuations" at 13.39pt}

% Korean section-heading font (overrides the manual's Latin \sectionfont).
\lazyplainfont\sectionfont{="Noto Sans CJK KR Light" at 18pt}
